% Source: https://artofproblemsolving.com/community/c3h2024723_geometry
% Stable identity: aops-topic-2024723-q01
Label the sides cyclically as $s_1,\dots,s_5$ (also used for their lengths), lying on lines $\ell_1,\dots,\ell_5$. In a convex polygon two adjacent sides are never parallel, and one side cannot be parallel to two others (those two would be parallel to each other, and three sides of a convex polygon cannot be parallel). So the two parallel pairs are disjoint pairs of non-adjacent sides. The non-adjacent pairs in a $5$-cycle are $\{1,3\},\{2,4\},\{3,5\},\{4,1\},\{5,2\}$, and every two disjoint ones have the form $\{i,i+2\},\{i+1,i+3\}$. After relabelling cyclically we may assume $\ell_1\parallel\ell_3$ and $\ell_2\parallel\ell_4$.

Let the vertices be $P_1=\ell_5\cap\ell_1$, $P_2=\ell_1\cap\ell_2$, $P_3=\ell_2\cap\ell_3$, $P_4=\ell_3\cap\ell_4$, $P_5=\ell_4\cap\ell_5$, and let $K=\ell_4\cap\ell_1$ (these lines are not parallel, since $\ell_4\parallel\ell_2$). Then $P_2P_3P_4K$ is a parallelogram.

\begin{asy}
usepackage("amsmath");
size(6cm);
defaultpen(fontsize(10pt));
pair K=(0,0), P2=(4.2,0), P3=(5.4,3), P4=(1.2,3);
pair P1=(1.6,0), P5=0.42*P4;
fill(P1--P2--P3--P4--P5--cycle, paleblue);
draw(P1--P2--P3--P4--P5--cycle, linewidth(0.9));
draw(K--P1, dashed); draw(K--P5, dashed);
dot(K); dot(P1); dot(P2); dot(P3); dot(P4); dot(P5);
label("$K$", K, SW); label("$P_1$", P1, S); label("$P_2$", P2, SE);
label("$P_3$", P3, NE); label("$P_4$", P4, NW); label("$P_5$", P5, W);
label("$s_1$", (P1+P2)/2, S); label("$s_2$", (P2+P3)/2, E); label("$s_3$", (P3+P4)/2, N);
label("$s_4$", (P4+P5)/2, W); label("$s_5$", (P5+P1)/2, NE);
label("$\ell_1$", (5.2,0), E); draw((4.2,0)--(5.1,0), gray);
label("\mbox{}",(min(currentpicture,user=true)+max(currentpicture,user=true))/2,black); // keeps the text after the figure black when LaTeX typesets the labels inline
\end{asy}

A convex polygon lies in the closed half-plane of each side line that contains it, so the pentagon lies in the strip between $\ell_1,\ell_3$ and in the strip between $\ell_2,\ell_4$, that is, in the parallelogram. Hence $P_5$, which is on $\ell_4$, lies on segment $P_4K$, and $P_1$ lies on segment $KP_2$. Also $P_5\ne K$: otherwise $P_5$ would lie on $\ell_1$ and $\ell_5$, forcing $P_5=\ell_1\cap\ell_5=P_1$. Similarly $P_1\ne K$. So $KP_5P_1$ is a genuine triangle, and the strict triangle inequality gives
\[
s_5=P_5P_1<KP_5+KP_1.
\]
Opposite sides of the parallelogram are equal, so $KP_4=P_2P_3=s_2$ and $KP_2=P_3P_4=s_3$. Therefore $KP_5=s_2-s_4$ and $KP_1=s_3-s_1$, and
\[
s_1+s_4+s_5<s_2+s_3.
\]

For the side lengths $2,2,2,3,3$, any three sides have total length at least $2+2+2=6$, while any two sides have total length at most $3+3=6$. So $s_1+s_4+s_5\ge6\ge s_2+s_3$, contradicting the inequality above. Hence no convex pentagon with sides $2,2,2,3,3$ has two pairs of parallel sides. In particular the ``maximum area of such pentagons'' asked about in the original post does not exist.
